\documentclass[11pt,a4paper]{article}
\usepackage[utf8]{inputenc}
\usepackage[T1]{fontenc}
\usepackage[catalan]{babel}
\usepackage{lmodern}
\usepackage[sfdefault,lining]{FiraSans}
\usepackage[a4paper,top=2.3cm,bottom=2.2cm,left=1.8cm,right=1.8cm,headheight=1.3cm,headsep=0.4cm,footskip=1cm]{geometry}
\usepackage{amsmath,amssymb}
\usepackage{icomma}
\usepackage{graphicx}
\usepackage[table]{xcolor}
\usepackage{tikz}
\usepackage{fontawesome5}
\usepackage{ccicons}
\usepackage{fancyhdr}
\usepackage{lastpage}
\usepackage{tasks}
\usepackage{needspace}
\usepackage[most]{tcolorbox}
\usepackage{array}
\usepackage[hidelinks,pdftitle={Fitxa d'exercicis de la unitat 1: Nombres reals. 1r de Batxillerat CCSS},pdfsubject={Matemàtiques Aplicades a les CCSS I},pdfkeywords={nombres reals, fracció generatriu, intervals, semirectes, valor absolut, unió i intersecció, radicals, forma exponencial, índex comú, racionalització, logaritmes}]{hyperref}

\definecolor{navy}{HTML}{1B2748}
\definecolor{accent}{HTML}{E8603A}
\definecolor{teal}{HTML}{1F9E90}
\definecolor{soft}{HTML}{EEF1F8}
\definecolor{muted}{HTML}{5D6685}

\setlength{\parindent}{0pt}
\setlength{\parskip}{3pt}
\settasks{label=\alph*), label-format=\bfseries\color{navy}, label-width=1.3em, item-indent=1.8em, column-sep=0.6em, after-item-skip=5pt, before-skip=3pt, after-skip=4pt}

% ---------- capçalera i peu ----------
\newcommand{\llicencia}{\href{https://creativecommons.org/licenses/by-nc-sa/4.0/deed.ca}{\ccbyncsa}\ \ Llicència Creative Commons Reconeixement-NoComercial-CompartirIgual 4.0 Internacional (CC BY-NC-SA 4.0)}
\pagestyle{fancy}
\fancyhf{}
\fancyhead[L]{\raisebox{-0.35\height}{\includegraphics[height=0.95cm]{logo-gva.png}}}
\fancyhead[C]{\small\color{muted}Fitxa d'exercicis · Unitat 1: Nombres reals · 1r Batxillerat CCSS}
\fancyhead[R]{\raisebox{-0.35\height}{\includegraphics[height=0.95cm]{logo-vermellar.png}}}
\fancyfoot[L]{\scriptsize\color{muted}\llicencia}
\fancyfoot[R]{\small\color{muted}\thepage}
\renewcommand{\headrulewidth}{0.4pt}
\renewcommand{\headrule}{\color{muted!40}\hrule height \headrulewidth}
\fancypagestyle{primera}{\fancyhf{}\renewcommand{\headrulewidth}{0pt}
  \fancyfoot[L]{\scriptsize\color{muted}\llicencia}
  \fancyfoot[R]{\small\color{muted}\thepage}}

% ---------- nivells i exercicis ----------
\newcommand{\estrelles}[1]{{\color{accent}\ifcase#1\or\faStar\faStar[regular]\faStar[regular]\or\faStar\faStar\faStar[regular]\or\faStar\faStar\faStar\fi}}
\newcounter{exer}
\newcommand{\ex}[1]{\par\needspace{4\baselineskip}\medskip\refstepcounter{exer}\noindent{\color{navy}\bfseries\theexer.}\ \ }

\newcommand{\seccio}[2]{\par\needspace{16\baselineskip}\bigskip
  \begin{tcolorbox}[colback=navy,colframe=navy,coltext=white,arc=2mm,boxrule=0pt,left=3mm,right=3mm,top=1.2mm,bottom=1.2mm]
  {\large\bfseries #1\quad #2}
  \end{tcolorbox}\vspace{-2pt}}

\newcommand{\exf}[1]{\par\medskip\refstepcounter{exer}\noindent{\color{navy}\bfseries\theexer.}\ \ }
\newcolumntype{C}[1]{>{\centering\arraybackslash}p{#1}}
% quadre de recordatori (només a la fitxa d'exercicis)
\newcommand{\recorda}[1]{\par\smallskip\begin{tcolorbox}[colback=accent!7,colframe=accent!55,arc=2mm,boxrule=0.5pt,left=2.5mm,right=2.5mm,top=1mm,bottom=1mm]\small{\color{accent}\faLightbulb[regular]}\ \textbf{Recorda:}\ #1\end{tcolorbox}}
\newcommand{\res}[1]{#1}\newcommand{\rest}[1]{#1}

\def\fmneed{11}% que el títol d'una secció no quede sol al final del full
% =====================================================================
% Capçalera FentMates de les fitxes: un color per a cada curs.
% Els colors són els de la pàgina inicial (LEVELS a assets/js/materials.js);
% el logo ve de _eines/logo/pdf/ (python3 _eines/logo/genera.py).
%
% Ús, al final del preàmbul de la unitat:
%   % =====================================================================
% Capçalera FentMates de les fitxes: un color per a cada curs.
% Els colors són els de la pàgina inicial (LEVELS a assets/js/materials.js);
% el logo ve de _eines/logo/pdf/ (python3 _eines/logo/genera.py).
%
% Ús, al final del preàmbul de la unitat:
%   % =====================================================================
% Capçalera FentMates de les fitxes: un color per a cada curs.
% Els colors són els de la pàgina inicial (LEVELS a assets/js/materials.js);
% el logo ve de _eines/logo/pdf/ (python3 _eines/logo/genera.py).
%
% Ús, al final del preàmbul de la unitat:
%   \input{../comu/fentmates.tex}
%   \fmcurs{1eso}
% i, al començament del document:
%   \capcalera{Matemàtiques}{1r ESO}{Fitxa d'exercicis · Unitat 2}{Divisibilitat}
% Opcional abans de \input: \def\fmneed{14} (espai mínim abans de cada secció).
% =====================================================================
\makeatletter
\def\fm@logodir{../../logo/pdf}
\providecommand{\fmneed}{10}
\def\fm@cursnom{}

% clau del curs, color i color del text damunt del color
\newcommand{\fm@curs}[3]{\@namedef{fm@color@#1}{#2}\@namedef{fm@text@#1}{#3}}
\fm@curs{1eso}{E6B84A}{navy}
\fm@curs{2eso}{3B82F6}{white}
\fm@curs{3eso}{FF7A50}{white}
\fm@curs{4eso-a}{A78BFA}{white}
\fm@curs{4eso-b}{6D28D9}{white}
\fm@curs{1batx-mi}{9B1C1C}{white}
\fm@curs{1batx-ccss}{EF6B6B}{white}
\fm@curs{2batx-mii}{1E6B47}{white}
\fm@curs{2batx-ccss}{2BB3A3}{white}

\newcommand{\fmcurs}[1]{%
  \@ifundefined{fm@color@#1}{\PackageError{fentmates}{Curs desconegut: #1}{}}{}%
  \edef\fm@tmp{\noexpand\definecolor{curs}{HTML}{\@nameuse{fm@color@#1}}%
               \noexpand\colorlet{curstext}{\@nameuse{fm@text@#1}}}\fm@tmp
  \colorlet{cursfosc}{curs!72!black}%
  \def\fm@logo{\fm@logodir/vertical-#1.pdf}%
  \def\fm@marca{\fm@logodir/marca-#1.pdf}}

% ---------- capçalera de la primera pàgina ----------
% #1 matèria (petit), #2 curs, #3 tipus de document i unitat, #4 títol
\newcommand{\capcalera}[4]{%
  \gdef\fm@cursnom{#2}%
  \par\noindent
  \begin{tcolorbox}[colback=curs,colframe=curs,boxrule=0pt,arc=2.2mm,width=3.5cm,height=2.15cm,
      valign=center,left=3mm,right=2mm,top=1mm,bottom=1mm,nobeforeafter,box align=center]
    \color{curstext}{\footnotesize #1}\par\vspace{1pt}{\fontsize{19}{21}\selectfont\bfseries #2\par}%
  \end{tcolorbox}\hspace{5mm}%
  \begin{minipage}[c]{\dimexpr\linewidth-3.5cm-5mm-2.9cm\relax}
    {\color{cursfosc}\small\bfseries\MakeUppercase{#3}}\par\vspace{2pt}
    {\color{navy}\fontsize{20}{23}\selectfont\bfseries #4\par}
  \end{minipage}\hfill
  \begin{minipage}[c]{2.6cm}\raggedleft\includegraphics[width=2.6cm]{\fm@logo}\end{minipage}\par
  \vspace{2pt}}

% ---------- barres de secció del color del curs ----------
\renewcommand{\seccio}[2]{\par\needspace{\fmneed\baselineskip}\bigskip
  \begin{tcolorbox}[colback=curs,colframe=curs,coltext=curstext,arc=2mm,boxrule=0pt,left=3mm,right=3mm,top=1.2mm,bottom=1.2mm]
  {\large\bfseries #1\quad #2}
  \end{tcolorbox}\vspace{-2pt}}

% ---------- estrelles de dificultat del color del curs (a Batxillerat no se'n posen) ----------
\@ifundefined{estrelles}{}{%
  \renewcommand{\estrelles}[1]{{\color{curs}\ifcase#1\or\faStar\faStar[regular]\faStar[regular]\or\faStar\faStar\faStar[regular]\or\faStar\faStar\faStar\fi}}}

% ---------- capçalera i peu de les pàgines ----------
% primera pàgina: logos institucionals al peu
\fancypagestyle{primera}{\fancyhf{}\renewcommand{\headrulewidth}{0pt}
  \fancyfoot[L]{\raisebox{-0.32\height}{\includegraphics[height=0.78cm]{logo-gva.png}}\hspace{5mm}%
    \raisebox{-0.32\height}{\includegraphics[height=0.78cm]{logo-vermellar.png}}}
  \fancyfoot[R]{\scriptsize\color{muted}\href{https://creativecommons.org/licenses/by-nc-sa/4.0/deed.ca}{\ccbyncsa}\ \ CC BY-NC-SA 4.0 · fentmates.org\qquad\small\thepage}}
% la resta: el símbol a l'esquerra i el curs a la dreta, amb el títol al centre com abans
\fancyhead[L]{\raisebox{-0.3\height}{\includegraphics[height=0.62cm]{\fm@marca}}}
\fancyhead[R]{\small\bfseries\color{cursfosc}\fm@cursnom}
\renewcommand{\headrule}{\color{curs!45}\hrule height \headrulewidth}
\makeatother

%   \fmcurs{1eso}
% i, al començament del document:
%   \capcalera{Matemàtiques}{1r ESO}{Fitxa d'exercicis · Unitat 2}{Divisibilitat}
% Opcional abans de \input: \def\fmneed{14} (espai mínim abans de cada secció).
% =====================================================================
\makeatletter
\def\fm@logodir{../../logo/pdf}
\providecommand{\fmneed}{10}
\def\fm@cursnom{}

% clau del curs, color i color del text damunt del color
\newcommand{\fm@curs}[3]{\@namedef{fm@color@#1}{#2}\@namedef{fm@text@#1}{#3}}
\fm@curs{1eso}{E6B84A}{navy}
\fm@curs{2eso}{3B82F6}{white}
\fm@curs{3eso}{FF7A50}{white}
\fm@curs{4eso-a}{A78BFA}{white}
\fm@curs{4eso-b}{6D28D9}{white}
\fm@curs{1batx-mi}{9B1C1C}{white}
\fm@curs{1batx-ccss}{EF6B6B}{white}
\fm@curs{2batx-mii}{1E6B47}{white}
\fm@curs{2batx-ccss}{2BB3A3}{white}

\newcommand{\fmcurs}[1]{%
  \@ifundefined{fm@color@#1}{\PackageError{fentmates}{Curs desconegut: #1}{}}{}%
  \edef\fm@tmp{\noexpand\definecolor{curs}{HTML}{\@nameuse{fm@color@#1}}%
               \noexpand\colorlet{curstext}{\@nameuse{fm@text@#1}}}\fm@tmp
  \colorlet{cursfosc}{curs!72!black}%
  \def\fm@logo{\fm@logodir/vertical-#1.pdf}%
  \def\fm@marca{\fm@logodir/marca-#1.pdf}}

% ---------- capçalera de la primera pàgina ----------
% #1 matèria (petit), #2 curs, #3 tipus de document i unitat, #4 títol
\newcommand{\capcalera}[4]{%
  \gdef\fm@cursnom{#2}%
  \par\noindent
  \begin{tcolorbox}[colback=curs,colframe=curs,boxrule=0pt,arc=2.2mm,width=3.5cm,height=2.15cm,
      valign=center,left=3mm,right=2mm,top=1mm,bottom=1mm,nobeforeafter,box align=center]
    \color{curstext}{\footnotesize #1}\par\vspace{1pt}{\fontsize{19}{21}\selectfont\bfseries #2\par}%
  \end{tcolorbox}\hspace{5mm}%
  \begin{minipage}[c]{\dimexpr\linewidth-3.5cm-5mm-2.9cm\relax}
    {\color{cursfosc}\small\bfseries\MakeUppercase{#3}}\par\vspace{2pt}
    {\color{navy}\fontsize{20}{23}\selectfont\bfseries #4\par}
  \end{minipage}\hfill
  \begin{minipage}[c]{2.6cm}\raggedleft\includegraphics[width=2.6cm]{\fm@logo}\end{minipage}\par
  \vspace{2pt}}

% ---------- barres de secció del color del curs ----------
\renewcommand{\seccio}[2]{\par\needspace{\fmneed\baselineskip}\bigskip
  \begin{tcolorbox}[colback=curs,colframe=curs,coltext=curstext,arc=2mm,boxrule=0pt,left=3mm,right=3mm,top=1.2mm,bottom=1.2mm]
  {\large\bfseries #1\quad #2}
  \end{tcolorbox}\vspace{-2pt}}

% ---------- estrelles de dificultat del color del curs (a Batxillerat no se'n posen) ----------
\@ifundefined{estrelles}{}{%
  \renewcommand{\estrelles}[1]{{\color{curs}\ifcase#1\or\faStar\faStar[regular]\faStar[regular]\or\faStar\faStar\faStar[regular]\or\faStar\faStar\faStar\fi}}}

% ---------- capçalera i peu de les pàgines ----------
% primera pàgina: logos institucionals al peu
\fancypagestyle{primera}{\fancyhf{}\renewcommand{\headrulewidth}{0pt}
  \fancyfoot[L]{\raisebox{-0.32\height}{\includegraphics[height=0.78cm]{logo-gva.png}}\hspace{5mm}%
    \raisebox{-0.32\height}{\includegraphics[height=0.78cm]{logo-vermellar.png}}}
  \fancyfoot[R]{\scriptsize\color{muted}\href{https://creativecommons.org/licenses/by-nc-sa/4.0/deed.ca}{\ccbyncsa}\ \ CC BY-NC-SA 4.0 · fentmates.org\qquad\small\thepage}}
% la resta: el símbol a l'esquerra i el curs a la dreta, amb el títol al centre com abans
\fancyhead[L]{\raisebox{-0.3\height}{\includegraphics[height=0.62cm]{\fm@marca}}}
\fancyhead[R]{\small\bfseries\color{cursfosc}\fm@cursnom}
\renewcommand{\headrule}{\color{curs!45}\hrule height \headrulewidth}
\makeatother

%   \fmcurs{1eso}
% i, al començament del document:
%   \capcalera{Matemàtiques}{1r ESO}{Fitxa d'exercicis · Unitat 2}{Divisibilitat}
% Opcional abans de \input: \def\fmneed{14} (espai mínim abans de cada secció).
% =====================================================================
\makeatletter
\def\fm@logodir{../../logo/pdf}
\providecommand{\fmneed}{10}
\def\fm@cursnom{}

% clau del curs, color i color del text damunt del color
\newcommand{\fm@curs}[3]{\@namedef{fm@color@#1}{#2}\@namedef{fm@text@#1}{#3}}
\fm@curs{1eso}{E6B84A}{navy}
\fm@curs{2eso}{3B82F6}{white}
\fm@curs{3eso}{FF7A50}{white}
\fm@curs{4eso-a}{A78BFA}{white}
\fm@curs{4eso-b}{6D28D9}{white}
\fm@curs{1batx-mi}{9B1C1C}{white}
\fm@curs{1batx-ccss}{EF6B6B}{white}
\fm@curs{2batx-mii}{1E6B47}{white}
\fm@curs{2batx-ccss}{2BB3A3}{white}

\newcommand{\fmcurs}[1]{%
  \@ifundefined{fm@color@#1}{\PackageError{fentmates}{Curs desconegut: #1}{}}{}%
  \edef\fm@tmp{\noexpand\definecolor{curs}{HTML}{\@nameuse{fm@color@#1}}%
               \noexpand\colorlet{curstext}{\@nameuse{fm@text@#1}}}\fm@tmp
  \colorlet{cursfosc}{curs!72!black}%
  \def\fm@logo{\fm@logodir/vertical-#1.pdf}%
  \def\fm@marca{\fm@logodir/marca-#1.pdf}}

% ---------- capçalera de la primera pàgina ----------
% #1 matèria (petit), #2 curs, #3 tipus de document i unitat, #4 títol
\newcommand{\capcalera}[4]{%
  \gdef\fm@cursnom{#2}%
  \par\noindent
  \begin{tcolorbox}[colback=curs,colframe=curs,boxrule=0pt,arc=2.2mm,width=3.5cm,height=2.15cm,
      valign=center,left=3mm,right=2mm,top=1mm,bottom=1mm,nobeforeafter,box align=center]
    \color{curstext}{\footnotesize #1}\par\vspace{1pt}{\fontsize{19}{21}\selectfont\bfseries #2\par}%
  \end{tcolorbox}\hspace{5mm}%
  \begin{minipage}[c]{\dimexpr\linewidth-3.5cm-5mm-2.9cm\relax}
    {\color{cursfosc}\small\bfseries\MakeUppercase{#3}}\par\vspace{2pt}
    {\color{navy}\fontsize{20}{23}\selectfont\bfseries #4\par}
  \end{minipage}\hfill
  \begin{minipage}[c]{2.6cm}\raggedleft\includegraphics[width=2.6cm]{\fm@logo}\end{minipage}\par
  \vspace{2pt}}

% ---------- barres de secció del color del curs ----------
\renewcommand{\seccio}[2]{\par\needspace{\fmneed\baselineskip}\bigskip
  \begin{tcolorbox}[colback=curs,colframe=curs,coltext=curstext,arc=2mm,boxrule=0pt,left=3mm,right=3mm,top=1.2mm,bottom=1.2mm]
  {\large\bfseries #1\quad #2}
  \end{tcolorbox}\vspace{-2pt}}

% ---------- estrelles de dificultat del color del curs (a Batxillerat no se'n posen) ----------
\@ifundefined{estrelles}{}{%
  \renewcommand{\estrelles}[1]{{\color{curs}\ifcase#1\or\faStar\faStar[regular]\faStar[regular]\or\faStar\faStar\faStar[regular]\or\faStar\faStar\faStar\fi}}}

% ---------- capçalera i peu de les pàgines ----------
% primera pàgina: logos institucionals al peu
\fancypagestyle{primera}{\fancyhf{}\renewcommand{\headrulewidth}{0pt}
  \fancyfoot[L]{\raisebox{-0.32\height}{\includegraphics[height=0.78cm]{logo-gva.png}}\hspace{5mm}%
    \raisebox{-0.32\height}{\includegraphics[height=0.78cm]{logo-vermellar.png}}}
  \fancyfoot[R]{\scriptsize\color{muted}\href{https://creativecommons.org/licenses/by-nc-sa/4.0/deed.ca}{\ccbyncsa}\ \ CC BY-NC-SA 4.0 · fentmates.org\qquad\small\thepage}}
% la resta: el símbol a l'esquerra i el curs a la dreta, amb el títol al centre com abans
\fancyhead[L]{\raisebox{-0.3\height}{\includegraphics[height=0.62cm]{\fm@marca}}}
\fancyhead[R]{\small\bfseries\color{cursfosc}\fm@cursnom}
\renewcommand{\headrule}{\color{curs!45}\hrule height \headrulewidth}
\makeatother

\fmcurs{1batx-ccss}

\begin{document}
\thispagestyle{primera}

\capcalera{Matemàtiques CCSS}{1r Batx}{Fitxa d'exercicis · Unitat 1}{\raggedright Nombres reals}
\vspace{10pt}
Nom i cognoms: \rule{8.4cm}{0.4pt}\hfill Grup: \rule{1.6cm}{0.4pt}\hfill Data: \rule{2.2cm}{0.4pt}
\vspace{8pt}

\begin{tcolorbox}[colback=soft,colframe=soft,arc=2mm,left=3mm,right=3mm,top=1.5mm,bottom=1.5mm]
\small Resol els exercicis al quadern mostrant tots els passos. Dona els resultats de manera exacta (fraccions i arrels simplificades, sense arrels al denominador) i, si cal, també aproximats. Escriu els conjunts com a intervals i representa'ls a la recta real.
\end{tcolorbox}

% =====================================================
\seccio{1}{Nombres reals}

\exf{1} Indica el menor dels conjunts numèrics ($\mathbb N$, $\mathbb Z$, $\mathbb Q$ o $\mathbb I$) als quals pertanyen els nombres següents:\begin{tasks}(5)
  \task $\sqrt5$
  \task $2$
  \task $-3$
  \task $24/6$
  \task $2,45$
  \task $-1,33$
  \task $2,3\widehat{45}$
  \task $\pi$
  \task $5/12$
  \task $\sqrt{\frac94}$
\end{tasks}

\ex{1} Raona si les afirmacions següents són verdaderes o falses:\begin{tasks}(1)
  \task Tot nombre enter és racional.
  \task Tot nombre decimal es pot expressar en forma de fracció.
  \task Tot nombre natural es pot escriure com una fracció.
  \task Si es sumen dos nombres racionals, el resultat és sempre un nombre racional.
  \task Si es multipliquen dos nombres irracionals, el resultat és sempre irracional.
\end{tasks}

\ex{1} Troba la fracció generatriu i simplifica-la:\begin{tasks}(3)
  \task $0,1\widehat6$
  \task $2,\widehat{09}$
  \task $-1,2\widehat{45}$
  \task $3,\widehat{12}$
  \task $0,4\widehat5$
  \task $2,45$
\end{tasks}\recorda{Fracció generatriu: al numerador, el nombre sense coma (fins al final del primer període) menys la part no periòdica; al denominador, tants 9 com xifres té el període i tants 0 com xifres té l'anteperíode.}

\ex{2} Realitza les operacions següents escrivint prèviament els nombres en forma de fracció:\begin{tasks}(2)
  \task $2,\widehat7+3,\widehat4-1,\widehat1$
  \task $1,\widehat8+0,\widehat1$
  \task $0,\widehat3\cdot1,\widehat5-0,1\widehat6$
  \task $\left(1,\widehat2+0,5\right)\cdot0,\widehat{18}$
\end{tasks}

\ex{2} Classifica cada nombre en el menor conjunt numèric. Opera primer!\begin{tasks}(3)
  \task $\sqrt2\cdot\sqrt8$
  \task $\sqrt2+\sqrt8$
  \task $\pi-\pi$
  \task $0,1010010001\ldots$
  \task $-\sqrt{\frac94}$
  \task $\sqrt{-4}$
\end{tasks}


% =====================================================
\seccio{2}{Intervals i semirectes}

\exf{1} Completa la taula següent:

\begin{center}\renewcommand{\arraystretch}{1.5}\begin{tabular}{|C{4.6cm}|C{3.2cm}|C{6cm}|}\hline\rowcolor{soft}\textbf{Desigualtat} & \textbf{Interval} & \textbf{Representació}\\\hline$\{x\in\mathbb R\ /\ x<-3\}$ & & \\\hline & $x\in[-2,\, 4)$ & \\\hline $\{x\in\mathbb R\ /\ -5<x\le-3\}$ & & \\\hline & $x\in(4,\, +\infty)$ & \\\hline\end{tabular}\end{center}

\ex{1} Escriu com a interval o semirecta i representa els conjunts numèrics següents:\begin{tasks}(2)
  \task Nombres majors que 3.
  \task $\{x\in\mathbb R/2\le x\le5\}$
  \task $\{x\in\mathbb R/3<x\le7\}$
  \task Nombres menors que 1 excloent-ne el 0.
  \task Nombres reals positius.
  \task $\{x\in\mathbb R/-2\le x<5\}$
  \task Nombres no negatius.
  \task $\{x\in\mathbb R/x\le-1\}$
\end{tasks}

\ex{1} Expressa com a desigualtat, com a interval i representa:\begin{tasks}(2)
  \task $x$ és menor que $-5$.
  \task 3 és menor o igual que $x$.
  \task $x$ està comprès entre $-5$ i 1.
  \task $x$ està entre $-2$ i 0, tots dos inclosos.
  \task $x$ és major o igual que $-1$ i menor que 6.
  \task $x$ no és major que 2.
\end{tasks}

\ex{1} Escriu la desigualtat que verifica tot nombre $x$ que pertany a aquests intervals i semirectes, i representa'ls:\begin{tasks}(4)
  \task $[-2,\, 7]$
  \task $[13,\, +\infty)$
  \task $(-\infty,\, 0)$
  \task $(-3,\, 0]$
  \task $\left[\frac32,\, 6\right)$
  \task $(-\infty,\, +\infty)$
  \task $(1,\, 6)$
  \task $(0,\, 4]$
  \task $(-\infty,\, 2)$
  \task $(-\infty,\, {-}5]$
  \task $(-1,\, +\infty)$
  \task $\left(2,\, \frac72\right)$
\end{tasks}

\ex{1} Quins nombres enters pertanyen a cadascun dels conjunts següents?\begin{tasks}(4)
  \task $[-2,\, 3)$
  \task $\left(-1,\, \frac52\right]$
  \task $(-4,\, {-}1)$
  \task $(3,\, 4)$
\end{tasks}


% =====================================================
\seccio{3}{Valor absolut}

\exf{1} Calcula:\begin{tasks}(4)
  \task $|-7|$
  \task $|3-8|$
  \task $-|-5|$
  \task $|-2|+|5|$
  \task $|(-3)\cdot4|$
  \task $\big||-5|-8\big|$
  \task $|\pi-4|$
  \task $|1-\sqrt3|$
\end{tasks}

\ex{1} Calcula la distància entre els parells de nombres següents:\begin{tasks}(4)
  \task $-4$ i $6$
  \task $2$ i $7$
  \task $-3$ i $-9$
  \task $-\frac12$ i $\frac32$
\end{tasks}

\ex{2} Resol les equacions següents:\begin{tasks}(3)
  \task $|x|=5$
  \task $|x-3|=2$
  \task $|2x+1|=7$
  \task $|x+4|=-1$
  \task $|3x-6|=0$
  \task $|5-x|=1$
\end{tasks}\recorda{Si $r>0$, \ $|A|=r\iff A=r$ o $A=-r$.}

\ex{2} Esbrina per a quins valors de $x$ es compleixen les desigualtats següents. Escriu la solució com a interval i representa-la:\begin{tasks}(3)
  \task $|x|<3$
  \task $|x|\ge3$
  \task $|x-2|\le3$
  \task $|x-4|>2$
  \task $|2x-1|>1$
  \task $|-5x+3|<7$
  \task $|x+1|\le\frac12$
  \task $|-5x+25|\le-2$
  \task $\left|x+\frac34\right|>-1$
\end{tasks}\recorda{Si $r>0$: \ $|A|<r\iff -r<A<r$ \quad i \quad $|A|>r\iff A<-r$ o $A>r$. En multiplicar o dividir per un nombre negatiu, la desigualtat canvia de sentit.}

\ex{2} Expressa cada conjunt amb una desigualtat amb valor absolut (busca el centre i el radi):\begin{tasks}(4)
  \task $(1,\, 7)$
  \task $[-5,\, 1]$
  \task $(-\infty,\, {-}1)\cup(5,\, +\infty)$
  \task $\left[-\frac12,\, \frac12\right]$
\end{tasks}


% =====================================================
\seccio{4}{Unió i intersecció d'intervals}

\exf{2} Donats els intervals $A=[-3,\, 2]$, $B=[0,\, 5]$, $C=[2,\, +\infty)$, $D=(0,\, 10)$ i $E=\left(2,\, \frac72\right)$, calcula:\begin{tasks}(3)
  \task $A\cap B$
  \task $A\cap D$
  \task $A\cup B$
  \task $(A\cup B)\cap C$
  \task $A\cap C$
  \task $B\cup C$
  \task $C\cap D$
  \task $E\cap B$
  \task $A\cup E$
\end{tasks}\recorda{Dibuixa els intervals l'un damunt de l'altre en la mateixa recta. La unió ($\cup$) són els nombres que estan en algun dels dos; la intersecció ($\cap$), els que estan en tots dos alhora.}

\ex{2} Donats $A=(-\infty,\, {-}4]$, $B=[2,\, 5)$, $C=[-6,\, 1]$, $D=[0,\, +\infty)$ i $E=(-2,\, {-}1)$, calcula:\begin{tasks}(3)
  \task $A\cup B$
  \task $B\cap C$
  \task $D\cap C$
  \task $A\cup(B\cap D)$
  \task $B\cap(A\cup E)$
  \task $B\cap D\cap E$
  \task $(B\cup C)\cap A$
  \task $A\cup C$
  \task $C\cap E$
  \task $A\cap D$
\end{tasks}

\ex{1} Escriu en forma d'interval els nombres que verifiquen:\begin{tasks}(3)
  \task $x<3$ i $x\ge5$
  \task $x>0$ i $x<4$
  \task $x\le-1$ i $x>1$
  \task $x<3$ i $x\ge-2$
  \task $x<-1$ o $x\ge2$
  \task $x>0$ o $x<4$
\end{tasks}

\ex{2} Siga $A$ el conjunt dels nombres que compleixen $|x-1|<2$ i $B=[2,\, +\infty)$. Escriu $A$ com a interval i calcula $A\cup B$ i $A\cap B$. Representa'ls.


% =====================================================
\seccio{5}{Radicals}

\exf{1} Calcula el valor dels radicals següents, si existeixen:\begin{tasks}(4)
  \task $\sqrt{16}$
  \task $\sqrt[3]{-8}$
  \task $\sqrt[5]{32}$
  \task $\sqrt[3]{\frac{27}{8}}$
  \task $\sqrt[3]{0,064}$
  \task $\sqrt[4]{81}$
  \task $\sqrt[4]{-16}$
  \task $\sqrt[5]{-\frac{1}{32}}$
\end{tasks}

\ex{1} Escriu en forma de potència els radicals de la primera fila i en forma de radical les potències de la segona:\begin{tasks}(5)
  \task $\sqrt5$
  \task $\sqrt[3]{x^2}$
  \task $\sqrt[5]{a^3}$
  \task $\dfrac{1}{\sqrt2}$
  \task $\dfrac{1}{\sqrt[3]{a^2}}$
  \task $3^{\frac25}$
  \task $a^{\frac34}$
  \task $5^{-\frac12}$
  \task $x^{1,5}$
  \task $2^{\frac53}$
\end{tasks}\recorda{$\sqrt[n]{a^m}=a^{\frac mn}$}

\ex{1} Calcula sense calculadora passant a exponents fraccionaris:\begin{tasks}(3)
  \task $8^{\frac23}$
  \task $81^{0,75}$
  \task $\left(\dfrac{1}{32}\right)^{-\frac25}$
  \task $16^{\frac34}$
  \task $27^{-\frac13}$
  \task $25^{\frac32}$
\end{tasks}

\ex{1} Simplifica els radicals següents:\begin{tasks}(4)
  \task $\sqrt[6]{2^4}$
  \task $\sqrt[9]{a^6}$
  \task $\sqrt[15]{8a^3}$
  \task $\sqrt[8]{16a^4}$
  \task $\sqrt[6]{0,027}$
  \task $\sqrt[12]{x^8y^4}$
  \task $\sqrt[4]{25}$
  \task $\sqrt[6]{64}$
\end{tasks}

\ex{2} Redueix els radicals de cada apartat al mínim índex comú:\begin{tasks}(2)
  \task $\sqrt2,\ \sqrt[3]3$
  \task $\sqrt[3]{n^2},\ \sqrt[4]{m^3},\ \sqrt[6]{n^5}$
  \task $\sqrt a,\ \sqrt[3]{a^2},\ \sqrt[4]{a^3}$
  \task $\sqrt[3]{5xy^2},\ \sqrt{6x^3z}$
\end{tasks}

\ex{2} Ordena de menor a major reduint a índex comú:\begin{tasks}(3)
  \task $\sqrt3,\ \sqrt[3]5,\ \sqrt[4]{10}$
  \task $\sqrt2,\ \sqrt[3]3,\ \sqrt[6]6$
  \task $\sqrt[3]2,\ \sqrt[4]3$
\end{tasks}

\ex{2} Extrau tots els factors possibles de cadascun dels radicals següents:\begin{tasks}(5)
  \task $\sqrt8$
  \task $\sqrt{12}$
  \task $\sqrt[3]{16}$
  \task $\sqrt[3]{54}$
  \task $\sqrt[5]{64}$
  \task $\sqrt{\frac{27}{4}}$
  \task $3\sqrt{8a^3}$
  \task $\sqrt[3]{54a^4}$
  \task $\sqrt{72x^3y^4}$
  \task $\sqrt[4]{32x^6}$
\end{tasks}

\ex{2} Introdueix el coeficient dins del radical i simplifica:\begin{tasks}(3)
  \task $7\sqrt a$
  \task $2a\sqrt{3a}$
  \task $x\sqrt{\frac1x}$
  \task $\frac32\sqrt{\frac23}$
  \task $2\sqrt[3]3$
  \task $3\sqrt[3]2$
  \task $\frac23\sqrt[3]{\frac{81}{4}}$
  \task $a^2\sqrt a$
  \task $3+\sqrt7$
\end{tasks}

\ex{2} Realitza les sumes i restes següents:\begin{tasks}(2)
  \task $\sqrt8+\sqrt{18}+\sqrt{32}+\sqrt{50}$
  \task $\sqrt{28}+\sqrt{63}-\sqrt{175}$
  \task $3\sqrt2-3\sqrt8+3\sqrt{18}$
  \task $2a\sqrt3-\sqrt{27a^2}+a\sqrt{12}$
  \task $5\sqrt{20}-2\sqrt{45}+\sqrt5$
  \task $\sqrt{18}+5\sqrt{12}-2\sqrt2-\sqrt{243}$
  \task $\sqrt{\frac13}+\sqrt{27}$
  \task $2\sqrt[3]{16}-\sqrt[3]{54}+\sqrt[3]{250}$
\end{tasks}\recorda{Només es poden sumar radicals semblants (mateix índex i mateix radicand). Abans, extrau-ne els factors.}

\ex{2} Calcula i simplifica (redueix a índex comú quan calga):\begin{tasks}(4)
  \task $\sqrt5\cdot\sqrt6$
  \task $\sqrt{12}\cdot\sqrt3$
  \task $\sqrt[3]4\cdot\sqrt[3]{16}$
  \task $\sqrt5\cdot\sqrt[3]6$
  \task $\sqrt{20}:\sqrt5$
  \task $\sqrt[4]{9x^5y^3}:\sqrt[4]{3x^3y}$
  \task $\dfrac{\sqrt2}{\sqrt[3]2}$
  \task $\sqrt a\cdot\sqrt[3]a\cdot\sqrt[4]a$
\end{tasks}

\ex{2} Calcula les potències i escriu davall d'un sol radical:\begin{tasks}(3)
  \task $(\sqrt3)^4$
  \task $(\sqrt[3]2)^6$
  \task $(\sqrt2)^5$
  \task $\left(\sqrt[4]{a^3}\right)^2$
  \task $(2\sqrt3)^2$
  \task $\sqrt[3]{(27x^3y^6)^5}$
  \task $\sqrt{\sqrt8}$
  \task $\sqrt[3]{\sqrt{32}}$
  \task $\sqrt{2\sqrt2}$
\end{tasks}

\ex{2} Opera i simplifica:\begin{tasks}(3)
  \task $(\sqrt3+\sqrt2)^2$
  \task $(2\sqrt3-1)(2\sqrt3+1)$
  \task $(\sqrt5-2)^2$
  \task $\sqrt2\left(\sqrt8-\sqrt{18}\right)$
  \task $(1+\sqrt2)^2-\sqrt8$
\end{tasks}


% =====================================================
\seccio{6}{Racionalització}

\exf{2} Racionalitza el denominador de les expressions següents:\begin{tasks}(5)
  \task $\dfrac{5}{\sqrt5}$
  \task $\dfrac{2\sqrt6}{\sqrt2}$
  \task $\dfrac{3}{\sqrt5}$
  \task $\sqrt{\dfrac52}$
  \task $\dfrac{3}{2\sqrt[3]4}$
  \task $\dfrac{6}{\sqrt[3]3}$
  \task $\dfrac{3y}{x\sqrt y}$
  \task $\dfrac{2\sqrt{12}}{5\sqrt3}$
  \task $\dfrac{10}{\sqrt[4]8}$
\end{tasks}\recorda{Si al denominador hi ha una sola arrel, multipliquem dalt i baix pel que calga perquè l'arrel desaparega: $\dfrac{1}{\sqrt[n]{a^m}}=\dfrac{\sqrt[n]{a^{n-m}}}{a}$.}

\ex{2} Racionalitza el denominador de les expressions següents:\begin{tasks}(3)
  \task $\dfrac{1}{\sqrt2+\sqrt3}$
  \task $\dfrac{2}{\sqrt2-1}$
  \task $\dfrac{\sqrt2}{3-\sqrt2}$
  \task $\dfrac{1-\sqrt2}{1+\sqrt2}$
  \task $\dfrac{4}{\sqrt5-1}$
  \task $\dfrac{5}{2(\sqrt3-\sqrt2)}$
  \task $\dfrac{2+\sqrt2}{2-\sqrt2}$
  \task $\dfrac{\sqrt3}{\sqrt3+1}$
  \task $\dfrac{1}{\sqrt x-\sqrt y}$
\end{tasks}\recorda{Si al denominador hi ha una suma o una resta amb arrels, multipliquem pel conjugat: $(a+b)(a-b)=a^2-b^2$.}

\ex{2} Racionalitza, opera i simplifica:\begin{tasks}(3)
  \task $\dfrac{\sqrt3+1}{\sqrt3-1}-\dfrac{\sqrt3-1}{\sqrt3+1}$
  \task $\dfrac{1}{\sqrt2}+\dfrac{1}{\sqrt2+1}$
  \task $\dfrac{2}{\sqrt3}-\dfrac{1}{2+\sqrt3}$
\end{tasks}


% =====================================================
\seccio{7}{Logaritmes}

\exf{1} Aplicant la definició, calcula el valor dels logaritmes següents:\begin{tasks}(4)
  \task $\log_2 8$
  \task $\log_3 81$
  \task $\log 1000$
  \task $\log 0,0001$
  \task $\ln e^{5}$
  \task $\log_4 16$
  \task $\log_4 0,25$
  \task $\log_5 1$
\end{tasks}\recorda{$\log_a b=c\iff a^c=b$ \ ($a>0$, $a\neq1$, $b>0$). $\log$ és el logaritme en base 10 i $\ln$, en base $e$.}

\ex{1} Aplicant la definició, calcula el valor dels logaritmes següents:\begin{tasks}(4)
  \task $\log_3\sqrt{27}$
  \task $\log_7\frac{1}{49}$
  \task $\log_3 3\sqrt3$
  \task $\log_3 0,\widehat3$
  \task $\log_5\sqrt[3]{25}$
  \task $\log_9\sqrt[3]3$
  \task $\log_{\frac12}8$
  \task $\log_8 0,125$
\end{tasks}

\ex{2} Calcula, si és possible, el valor de $x$ en cadascuna de les expressions següents:\begin{tasks}(4)
  \task $\log_x 8=3$
  \task $\log_x 8=-3$
  \task $\log_5 x=-1$
  \task $\log_2 x=5$
  \task $\log_x 49=2$
  \task $\log_3(-81)=x$
  \task $\log_x\sqrt2=0$
  \task $\log x=-2$
\end{tasks}\recorda{$\log_a b$ només existeix si $a>0$, $a\neq1$ i $b>0$.}

\ex{1} Amb l'ajuda de la calculadora i el canvi de base, obtín aproximacions fins a les mil·lèsimes:\begin{tasks}(3)
  \task $\log_3 20$
  \task $\log_{0,5}60$
  \task $\log_{\sqrt2}3$
  \task $\log_7 100$
  \task $\log_2 0,3$
  \task $\log_5 2$
\end{tasks}\recorda{Canvi de base: $\log_a x=\dfrac{\log x}{\log a}=\dfrac{\ln x}{\ln a}$.}


% =====================================================
\seccio{8}{Propietats dels logaritmes}

\exf{1} Troba el resultat de les expressions següents mitjançant les propietats dels logaritmes:\begin{tasks}(2)
  \task $2\log_4 16+\log_2 32-3\log_7 49$
  \task $\log_2 8+\log_3 27+3\log_5 125$
  \task $\log_5 625-\log_9 81+\log_8 64$
  \task $\log_6 4+\log_6 9$
  \task $\log_2 40-\log_2 5$
  \task $\log 20+\log 5$
  \task $\log_3 54-\log_3 2$
  \task $2\log 5+\log 4$
\end{tasks}\recorda{$\log_a(p\cdot q)=\log_a p+\log_a q$, \ $\log_a\frac pq=\log_a p-\log_a q$, \ $\log_a p^n=n\log_a p$.}

\ex{2} Considerant $\log2=0,301$ i $\log3=0,477$, calcula sense calculadora:\begin{tasks}(5)
  \task $\log4$
  \task $\log6$
  \task $\log5$
  \task $\log12$
  \task $\log1,5$
  \task $\log0,02$
  \task $\log3,6$
  \task $\log\sqrt3$
  \task $\log0,\widehat3$
  \task $\log45$
\end{tasks}

\ex{2} Si $\log_3 2=0,631$ i $\log_3 5=1,465$, calcula $\log_3 150$ i $\log_3\frac{10}{9}$.

\ex{2} Pren logaritmes decimals i desenvolupa les expressions:\begin{tasks}(2)
  \task $P=10x^3y$
  \task $Q=\dfrac{100x^2}{x+y}$
  \task $R=\dfrac{x^2\sqrt y}{z^3}$
  \task $S=\sqrt[3]{a^2b}$
\end{tasks}

\ex{2} Expressa el valor de $E$ en cada cas sense que apareguen logaritmes:\begin{tasks}(1)
  \task $\log E=2-2\log x+\log y$
  \task $\log E=3\log2-4\log x+3\log y$
  \task $\log E=\log(x-2y)+\log(x+y)$
  \task $\log E=\frac12\log x+\log3$
  \task $\log E=1+\log a-2\log b$
\end{tasks}


% =====================================================
\seccio{9}{Problemes}

\exf{2} La temperatura $T$ d'una nevera ha de ser de $4\,^\circ$C amb un marge d'error de $1,5\,^\circ$C. Escriu aquesta condició amb un valor absolut i com a interval. És acceptable una lectura de $5,7\,^\circ$C?

\ex{2} Un quadrat té una àrea de $50\ \text{cm}^2$. Calcula'n el costat i la diagonal en forma de radical simplificat.

\ex{2} Ingressem 2\,000 € en un compte al 3\,\% d'interès compost anual. Quants anys han de passar, com a mínim, perquè hi haja 3\,000 €?

\ex{2} Un cotxe costa 20\,000 € i perd un 15\,\% del seu valor cada any. Quants anys han de passar perquè valga menys de 8\,000 €?

\ex{2} Una ciutat té 50\,000 habitants i la població creix un 2\,\% cada any. Quants anys tardarà a arribar a 60\,000 habitants?

\ex{2} Un bacteri es duplica cada hora si té les condicions òptimes. Quantes hores hauran de passar perquè un bacteri inicial cree una colònia de mil milions de bacteris?

\ex{1} Troba l'error en cadascun d'aquests càlculs i corregeix-lo:\begin{tasks}(2)
  \task $\sqrt{9+16}=3+4=7$
  \task $-3^2=9$
  \task $\log(5+5)=\log5+\log5$
  \task $2^{-3}=-8$
  \task $\sqrt2+\sqrt3=\sqrt5$
  \task $[1,\, 3)\cap(3,\, 5]=\{3\}$
\end{tasks}

\ex{3} Sense calculadora, ordena de menor a major $\log_2 10$, \ $3$ \ i \ $\log_3 10$. Justifica-ho.

\vfill
\begin{center}\color{muted}\small\faCheckCircle[regular]\ \ Ànim! Comprova els teus resultats amb les eines interactives dels apunts de la unitat.\end{center}

\end{document}
